\section{Related Work}

\subsection{From Categorical SER to Self-Supervised Representations}

Speech emotion recognition (SER) has evolved significantly from early systems relying on hand-crafted acoustic features \cite{elayadi2011survey, ververidis2006emotional} to end-to-end deep learning architectures that automatically learn discriminative representations \cite{trigeorgis2016adieu}. In recent years, self-supervised learning (SSL) has reshaped the field by enabling models to acquire general-purpose speech representations from vast quantities of unlabeled audio. Architectures such as wav2vec~2.0 \cite{Wav2vec2}, HuBERT \cite{HuBERT}, and WavLM \cite{WavLM} successfully leveraged masked prediction to mitigate the data scarcity inherent in supervised approaches. More recently, specialized models such as \textit{emotion2vec} \cite{emotion2vec} have introduced distillation frameworks explicitly targeting universal speech emotion representations, achieving robust benchmark performance. However, despite their acoustic grounding, these SSL approaches remain constrained by task-specific classification heads that map continuous human affect into rigid, discrete categories \cite{SUPERB}. This categorical bottleneck leaves limited room for the expressive paralinguistic context that matters in real interactions, such as vocal tension, breathiness, sarcasm, and dynamic tonal shifts.

\subsection{Speech Language Models and the Hallucination Bottleneck}

To address the limitations of discrete categorization, the emergence of multimodal large language models (MLLMs) and specialized speech language models (SLMs) represents a paradigm shift, reframing emotion inference as an instance of general auditory reasoning. Foundational models such as SALMONN \cite{SALMONN} and Qwen-Audio \cite{Qwen-Audio} established an encoder-connector-LLM architecture, enabling systems to generate free-form text descriptions of affective states. Subsequent scale expansions have pushed the boundaries of audio understanding; for instance, Step-Audio~2 \cite{Step-Audio2} incorporates a latent audio encoder for fine-grained paralinguistic recognition, Kimi-Audio \cite{Kimi-Audio} utilizes extensive flow-matching pre-training, and MOSS-Audio \cite{MOSS-Audio} introduces time-marker insertion to preserve temporal prosody. While these SLMs demonstrate the potential for rich, open-ended emotion description, they can still suffer from insufficient acoustic grounding and paralinguistic hallucinations, generating affective narratives that are not fully supported by the audio signal \cite{Qwen2-Audio}.

\subsection{Emotion Alignment and Naturalistic Data}

To refine SLM capabilities and reduce cross-modal hallucinations, researchers have increasingly explored Direct Preference Optimization (DPO) and targeted instruction tuning \cite{DPO}. DPO aligns model outputs with human judgments by contrasting preferred and rejected responses, a technique recently adapted for audio applications such as emotional text-to-speech \cite{gao2024emodpo}. Concurrently, the foundational data driving these models has shifted from small, acted corpora such as IEMOCAP \cite{IEMOCAP} to large-scale, naturalistic datasets such as MSP-Podcast \cite{MSP-Podcast}. These ``in-the-wild'' datasets expose models to the full complexity of real-world acoustic variability. Amphion EmoAnalyzer builds on these trends by combining descriptive emotion analysis, naturalistic bilingual data curation, and emotion-focused preference alignment into a unified model for more grounded and expressive speech emotion understanding.
